\chapter{Metodolog\'ia}
\label{cap:metodologia}

Este cap\'itulo integra, en formato de informe acad\'emico, el borrador
metodol\'ogico de Jaime (\texttt{docs/informe/borradores/jaime/metodologia-y-amenazas.md}),
apoyado en el checklist ya completado
\texttt{docs/checklists/ralph2021-engineering-research.md}~\cite{ralph2021-empirical-standards}.

\section{Enfoque metodol\'ogico}

BIOPET corresponde, de forma predominante, a una investigaci\'on de
\textbf{ingenier\'ia orientada al dise\~no y evaluaci\'on de un artefacto de
software} (\emph{Engineering Research}, tambi\'en denominada \emph{Design
Science} en la literatura de ingenier\'ia de software emp\'irica). El
equipo dise\~n\'o, construy\'o y evalu\'o emp\'iricamente un sistema real
--- backend Spring~Boot, frontend Angular e infraestructura Docker ---
frente a requisitos expl\'icitos (\texttt{docs/requisitos/SRS.md}), usando
m\'ultiples fuentes de evidencia t\'ecnica: pruebas automatizadas y
cobertura (JaCoCo), rendimiento (k6), usabilidad (SUS), calidad web
autom\'atica (Lighthouse), y seguridad (OWASP, ZAP, SpotBugs/Find Security
Bugs, auditor\'ia de SQL din\'amico).

\subsection*{Por qu\'e \emph{Engineering Research} y no una alternativa}
\begin{itemize}[nosep]
  \item \textbf{No es un experimento controlado.} Ninguna medici\'on
    manipula deliberadamente una variable independiente con asignaci\'on
    aleatoria a grupos de tratamiento/control. Las corridas de k6 miden el
    comportamiento del propio sistema bajo dos condiciones observacionales
    (cach\'e fr\'ia y caliente), no un dise\~no experimental con grupos
    comparados.
  \item \textbf{No es un estudio de caso observacional tradicional.} BIOPET
    es un artefacto construido por el propio equipo dentro de un Proyecto
    Fin de Curso, no un fen\'omeno observado \emph{in situ} en una
    cl\'inica veterinaria real ya operando con el sistema.
  \item \textbf{S\'i encaja con el patr\'on de \emph{Engineering
    Research}}, cuyos atributos esenciales (descripci\'on del artefacto,
    justificaci\'on de su relevancia, evaluaci\'on conceptual y emp\'irica,
    discusi\'on de limitaciones) est\'an todos presentes y verificados en
    el checklist Ralph et al. ya completado.
\end{itemize}

\subsection*{Afirmaciones que este informe evita deliberadamente}
No se afirma que BIOPET sea un experimento controlado (sin grupos de
control ni asignaci\'on aleatoria); no se afirma validaci\'on industrial
(toda la evaluaci\'on se realiz\'o en un entorno acad\'emico, sin usuarios
ni profesionales de una organizaci\'on externa real); no se afirma ninguna
forma de certificaci\'on (ISO/IEC~25010 se usa como marco de
clasificaci\'on de calidad, no como certificaci\'on obtenida); no se afirma
causalidad experimental en ninguna medici\'on; y no se afirma que los
resultados sean universalmente generalizables.

\section{Design Science Research: las seis actividades de Peffers et al.}
\label{sec:dsr-peffers}

\begin{table}[H]
  \centering
  \small
  \caption{Mapeo de BIOPET sobre las seis actividades de Peffers et al.~\cite{peffers2007dsrm}.}
  \label{tab:peffers}
  \begin{tabular}{@{}p{3.2cm}p{1.6cm}p{7.8cm}@{}}
    \toprule
    Actividad & Estado & Evidencia \\
    \midrule
    1. Identificaci\'on del problema y motivaci\'on & Completada & 15 observaciones docentes reales registradas y cerradas (\texttt{docs/observaciones/OBSERVACIONES.md}); SRS \S1 \\
    2. Definici\'on de objetivos de la soluci\'on & Completada & Requisitos REQ-F/REQ-NF verificables, trazados a HU/CU (\texttt{docs/trazabilidad/matriz.csv}) \\
    3. Dise\~no y desarrollo & Completada & Modelo C4 (contexto/contenedores/componentes), 7 ADR, c\'odigo fuente real en \texttt{Backend/src/}, \texttt{frontend/src/} \\
    4. Demostraci\'on & Completada & \texttt{make up} de punta a punta; colecciones Postman; k6 contra el backend real; ZAP contra el contenedor real \\
    5. Evaluaci\'on & Completada & JaCoCo, k6, SUS, Lighthouse, OWASP/ZAP/SpotBugs, todos con datos reales (secci\'on~\ref{sec:diseno-evaluacion}) \\
    6. Comunicaci\'on & Completada (este documento) & Este informe reemplaza al de la Tercera Entrega, con cifras actualizadas de v1.0.0 \\
    \bottomrule
  \end{tabular}
\end{table}

\section{Goal--Question--Metric (GQM)}
\label{sec:gqm}

\textbf{Goal:} analizar el artefacto BIOPET (backend Spring~Boot +
frontend Angular + infraestructura Docker) \textbf{con el prop\'osito de}
evaluar \textbf{con respecto a} su calidad funcional, mantenibilidad,
rendimiento, usabilidad, calidad web y seguridad, \textbf{desde el punto de
vista} del equipo de desarrollo y de evaluadores acad\'emicos,
\textbf{en el contexto de} un Proyecto Fin de Curso ejecutado en un
entorno de desarrollo local con Docker y desplegado en producci\'on sobre
el plan gratuito de Render~\cite{vansolingen2002gqm}.

\begin{table}[H]
  \centering
  \tiny
  \caption{Preguntas y m\'etricas GQM, con el valor de evidencia vigente al cierre de la Entrega Final.}
  \label{tab:gqm}
  \begin{tabular}{@{}p{2.0cm}p{4.6cm}p{4.6cm}@{}}
    \toprule
    Dimensi\'on & Pregunta & Valor de evidencia (fuente) \\
    \midrule
    Calidad funcional & \textquestiondown Los requisitos declarados tienen respaldo trazable? & Matriz consistente frente al SRS (\texttt{scripts/validate-traceability.sh}) \\
    Calidad funcional & \textquestiondown La suite de pruebas pasa consistentemente? & 205 pruebas, 0 fallos, 0 errores (\texttt{Backend/target/surefire-reports/}) \\
    Mantenibilidad & \textquestiondown Qu\'e proporci\'on del c\'odigo se ejercita? & LINE 91{,}80\,\% / BRANCH 79{,}39\,\%, umbral 70\,\% (\texttt{docs/mediciones/jacoco/METRICS.md}) \\
    Mantenibilidad & \textquestiondown El proceso de build es reproducible y automatizado? & CI con 6 jobs verdes (\texttt{.github/workflows/ci.yml}) \\
    Rendimiento & \textquestiondown Latencia y error del endpoint cacheado bajo carga? & p95 6{,}4--21{,}1~ms (caliente/fr\'io), error 0{,}0\,\% en 10 corridas (\texttt{docs/mediciones/perf/REPORT.md}) \\
    Usabilidad & \textquestiondown Qu\'e nivel de usabilidad percibida presenta el frontend seg\'un los registros de la evaluaci\'on SUS? & SUS media 74{,}44/100, IC95\,\% $[63{,}33,85{,}56]$, n=18 (\texttt{docs/mediciones/sus/REPORT.md}) \\
    Calidad web & \textquestiondown El frontend cumple umbrales base de calidad web? & Performance/Accessibility/Best Practices/SEO cumplidos en las 12 corridas mobile+desktop del 2026-08-18 (\texttt{docs/mediciones/lighthouse/}) \\
    Seguridad (din\'amica) & \textquestiondown Hay hallazgos de severidad alta por escaneo din\'amico? & 0 alertas altas, 1 informativa, ZAP Baseline (\texttt{docs/mediciones/sec/zap/README.md}) \\
    Seguridad (est\'atica) & \textquestiondown El an\'alisis est\'atico detecta patrones de inyecci\'on SQL? & 66 hallazgos totales, 0 de tipo \texttt{SQL\_*} (\rutalarga{docs/mediciones/sec/static-analysis/README.md}) \\
    Seguridad (SQL din\'amico) & \textquestiondown Los procedimientos almacenados construyen SQL de forma segura? & 0 hallazgos, \texttt{exit 0} (\texttt{scripts/audit-sql-dynamic.sh}) \\
    \bottomrule
  \end{tabular}
\end{table}

\section{Dise\~no de evaluaci\'on}
\label{sec:diseno-evaluacion}

Para cada fuente de evidencia emp\'irica se documenta prop\'osito, unidad
evaluada, procedimiento reproducible, m\'etrica principal y limitaci\'on
conocida (detalle ampliado por dimensi\'on en los
cap\'itulos~\ref{cap:pruebas-calidad} y \ref{cap:seguridad}):

\begin{enumerate}[nosep]
  \item \textbf{Pruebas automatizadas:} \texttt{mvn clean verify}
    (JUnit~5 + MockMvc + 2 pruebas de integraci\'on con Testcontainers
    contra PostgreSQL real). M\'etrica: pruebas/fallos/errores.
  \item \textbf{Cobertura JaCoCo:} regla \texttt{BUNDLE} sobre todo el
    m\'odulo backend; contadores LINE/BRANCH/COMPLEXITY.
  \item \textbf{Rendimiento (k6):} 5 corridas en fr\'io + 5 en caliente
    contra \texttt{GET /api/mascotas} sobre HTTPS/TLS~1.3; IC~95\,\% con
    distribuci\'on t de Student; prueba de Wilcoxon pareada entre
    condiciones fr\'ia/caliente.
  \item \textbf{Usabilidad (SUS):} cuestionario de Brooke (1996) sin
    modificar, aplicado a 18 participantes (P01--P18) que, seg\'un
    declaraci\'on del equipo, son externos al equipo; la disponibilidad
    de consentimiento informado individual firmado para cada
    participante no pudo confirmarse documentalmente durante una
    auditor\'ia posterior (ver secci\'on~\ref{sec:sus-final} y
    \texttt{docs/etica/regularizacion-sus/}).
  \item \textbf{Calidad web (Lighthouse):} \texttt{npx @lhci/cli autorun}
    contra el frontend servido por contenedor real, perfiles m\'ovil
    simulado y escritorio, 3 corridas por ruta y por perfil (12 corridas
    en la re-ejecuci\'on final del 2026-08-18).
  \item \textbf{Seguridad din\'amica (ZAP Baseline):} escaneo no
    autenticado contra el backend real dentro de la red Docker Compose.
  \item \textbf{Seguridad est\'atica (SpotBugs + Find Security Bugs):}
    an\'alisis de bytecode compilado, \texttt{effort=Max},
    \texttt{threshold=Low}.
  \item \textbf{Auditor\'ia de SQL din\'amico:} an\'alisis de patrones
    inseguros (\texttt{EXECUTE}, concatenaci\'on) sobre \texttt{db/procs/*.sql}.
  \item \textbf{Trazabilidad de requisitos:} validaci\'on cruzada
    autom\'atica SRS $\leftrightarrow$ matriz.
\end{enumerate}

Cada m\'etodo declara su propia limitaci\'on en el cap\'itulo
correspondiente y en el cap\'itulo~\ref{cap:amenazas-limitaciones-final}; ninguna
cifra de este informe se present\'o sin una fuente citada.

\subsection{Muestreo SUS: reclutamiento, poblaci\'on y l\'imites de generalizaci\'on}
\label{sec:muestreo-sus-metodologia}

Los 18 participantes (P01--P18) de la evaluaci\'on SUS
(secci\'on~\ref{sec:sus-final}) se reclutaron por \textbf{conveniencia,
dentro del c\'irculo cercano al equipo del proyecto}
(\texttt{docs/mediciones/sus/REPORT.md}, secci\'on ``Amenazas a la
validez''): no hubo aleatorizaci\'on, plataforma de reclutamiento pagada
ni acceso a usuarios de una cl\'inica veterinaria real. Se declara
expl\'icitamente hasta d\'onde llega esta evidencia, sin completar
detalles de canal de reclutamiento (por ejemplo, redes sociales
espec\'ificas o convocatoria formal) que no quedaron registrados y que el
equipo no puede sostener con evidencia.

\textbf{Poblaci\'on de origen (nivel de generalidad sostenible por la
evidencia):} adultos con acceso a un dispositivo con navegador web y
disposici\'on a colaborar con el equipo, sin relaci\'on documentada con
el dominio veterinario ni con la operaci\'on real de una cl\'inica. La
muestra result\'o demogr\'aficamente heterog\'enea dentro de ese c\'irculo
(\texttt{docs/mediciones/sus/REPORT.md}): edad 19--61 a\~nos (media
34{,}3), sexo balanceado (F=9, M=9), experiencia previa con aplicaciones
web variada (avanzada=6, intermedia=7, b\'asica=3, ninguna=2) y
dispositivo variado (laptop=10, escritorio=4, tablet=2, celular=2) --- lo
que \textbf{atenu\'a, sin eliminar}, la homogeneidad t\'ipica de una
muestra de conveniencia reclutada dentro de un c\'irculo social \'unico.

\textbf{Por qu\'e se us\'o muestreo por conveniencia y qu\'e sesgos
introduce.} Es el m\'etodo real y disponible para un equipo de tres
estudiantes en un Proyecto Fin de Curso, sin presupuesto para una
plataforma de reclutamiento ni acceso institucional a personal o
clientes de una cl\'inica veterinaria real; se declara como limitaci\'on
de recursos, no como decisi\'on metodol\'ogica \'optima. Introduce al
menos dos sesgos concretos: (1) \textbf{sesgo de selecci\'on} hacia
personas con alguna conexi\'on social al equipo, que pueden compartir
rasgos (nivel educativo, familiaridad general con tecnolog\'ia) no
representativos de la poblaci\'on general; y (2) \textbf{posible sesgo de
complacencia/deseabilidad social}, ya declarado en
\texttt{docs/mediciones/sus/REPORT.md}: participantes con relaci\'on
personal con el equipo pueden puntuar de forma m\'as favorable que un
usuario sin esa relaci\'on.

\textbf{A qu\'e poblaci\'on se pueden y no se pueden extrapolar los
resultados.} Los resultados (SUS medio 74{,}44/100, categor\'ia
``Bueno'') son razonablemente informativos sobre la percepci\'on de
usabilidad de \textbf{adultos con acceso a la web y diversidad et\'aria y
de experiencia digital similar a la observada en la muestra}, al
completar una tarea de \emph{onboarding} controlada (inicio de
sesi\'on, alta, edici\'on y baja l\'ogica de una mascota) en un entorno
sin presi\'on operativa real. \textbf{No se pueden extrapolar} a
usuarios reales de una cl\'inica veterinaria operando bajo carga de
trabajo real (veterinarios, auxiliares o personal administrativo
durante una jornada cl\'inica), ni a due\~nos de mascota en el momento de
una urgencia real, contextos que la evaluaci\'on no reprodujo; tampoco se
document\'o la ubicaci\'on geogr\'afica ni la relaci\'on de los
participantes con el sector veterinario, por lo que ese eje de
representatividad tampoco puede argumentarse en ning\'un sentido.

\textbf{Respaldo bibliogr\'afico.} Ralph, Baltes et al.
(2021)~\cite{ralph2021-empirical-standards}, en su est\'andar de
Ingenier\'ia de Software Emp\'irica, tratan el muestreo por conveniencia
como una pr\'actica leg\'itima y frecuente en investigaci\'on de
ingenier\'ia de software \textbf{siempre que la amenaza se reporte de
forma expl\'icita y la generalizaci\'on reclamada se acote al alcance que
la muestra realmente sostiene} --- exactamente el criterio aplicado en
esta secci\'on: no se afirma representatividad de la poblaci\'on general
ni del dominio veterinario, solo del segmento de poblaci\'on
efectivamente alcanzado.

\textbf{L\'imite \'etico declarado (sin evidencia inventada).} Como ya
se documenta en la secci\'on~\ref{sec:sus-final} y en
\texttt{docs/etica/regularizacion-sus/}, no existe evidencia verificable
de consentimiento informado individual firmado por cada participante, y
la constancia de regularizaci\'on firmada por el equipo \textbf{no
sustituye} esos consentimientos. Esta secci\'on no afirma una aprobaci\'on
\'etica previa ni consentimientos que no existan; el reclutamiento por
conveniencia descrito arriba no depende de ---ni corrige--- ese vac\'io
documental, que sigue declarado como limitaci\'on abierta.

\section{Reproducibilidad}

El sistema completo se reconstruye desde una clonaci\'on limpia con un
\'unico comando (\texttt{make up}), con im\'agenes de terceros fijadas por
\emph{digest} \texttt{sha256} para evitar derivas silenciosas entre
reconstrucciones. Los scripts que generan cada pieza de evidencia
(\rutalarga{scripts/archive-jacoco-evidence.sh},
\rutalarga{scripts/run-zap-baseline.sh}, \rutalarga{scripts/audit-sql-dynamic.sh},
\rutalarga{scripts/validate-traceability.sh}, \rutalarga{scripts/perf-analysis.py},
\rutalarga{scripts/analisis-sus.py}, \rutalarga{scripts/run-lighthouse.sh}, entre
otros) est\'an versionados, de modo que la evidencia no es solo consultable
sino regenerable. El detalle completo de reproducibilidad de despliegue se
documenta en el cap\'itulo~\ref{cap:despliegue}.

\section{Amenazas a la validez (resumen metodol\'ogico)}

Las amenazas a la validez de cada medici\'on se documentan en detalle en
el cap\'itulo~\ref{cap:amenazas-limitaciones-final}, clasificadas seg\'un
las cuatro categor\'ias est\'andar (interna, externa, de constructo y de
conclusi\'on). Este cap\'itulo de metodolog\'ia se limita a declarar el
principio que las gobierna: cada amenaza documentada incluye su efecto
posible, la mitigaci\'on ya aplicada (si existe) y el riesgo residual que
persiste, sin minimizar ninguna de ellas.
